\chapter{Keratosis Pilaris}

\section*{Definition and Overview}
Keratosis pilaris is a common, harmless skin condition that causes small, rough bumps, often described as chicken skin, most often on the upper arms, thighs, buttocks, and sometimes the cheeks. It occurs when keratin, a protein in the skin, blocks the openings of the hair follicles, forming tiny plugs. The bumps are usually painless and itch-free, though the skin can feel dry and rough. Keratosis pilaris is very common, often runs in families, and tends to improve with age. It cannot be cured, but it responds well to regular moisturising, gentle exfoliation, and creams containing urea or lactic acid.

\section*{Epidemiology}
Keratosis pilaris is one of the most common skin conditions, affecting up to 40 per cent of adults and around 50 to 80 per cent of adolescents. It is more common in people with dry skin, eczema, and asthma, and it often runs in families. It typically appears in childhood or the teenage years and may improve in adulthood, though some people have it lifelong. It affects all ethnic groups, and the appearance varies from barely visible to prominent and widespread. Most people never seek medical care for it.

\section*{Aetiology and Pathophysiology}
Keratosis pilaris results from abnormal keratinisation of the hair follicles: keratin, the structural protein that forms the skin's outer layer, builds up and blocks the follicle opening, creating a small plug and a bump. The exact cause is unknown, but it is related to genetics and to the same tendency to dry skin that underlies eczema. The condition is more noticeable in winter, when skin is drier, and with low humidity, and it may improve in summer. The follicles are not infected or inflamed in the usual sense, which is why the condition is harmless, and the plugs can be eased by exfoliation and moisturising.

\section*{Clinical Features}
\begin{itemize}
    \item Small, rough, red or skin-coloured bumps on the upper arms, thighs, buttocks, or cheeks
    \item A chicken-skin or goosebump texture
    \item Dry, rough skin over the affected areas
    \item Bumps that are usually painless and do not itch
    \item Redness around the bumps in some people
    \item Worse in winter and with dry skin
    \item Often associated with eczema, dry skin, or a family history
    \item The condition is harmless and does not scar
\end{itemize}

\section*{Differential Diagnosis}
Other conditions can cause bumps or rough skin.
\begin{itemize}
    \item \textbf{Acne:} inflamed pimples and blackheads
    \item \textbf{Folliculitis:} infected hair follicles with redness and pus
    \item \textbf{Eczema:} itchy, inflamed patches
    \item \textbf{Psoriasis:} scaly plaques
    \item \textbf{Allergic reactions:} hives and contact rashes
    \item \textbf{Other keratoses:} thickened skin conditions
    \item \textbf{Lichen planus and other papular rashes:} uncommon conditions
\end{itemize}

\section*{When to Seek Medical Care}
Keratosis pilaris does not usually need medical treatment, but a doctor or dermatologist can confirm the diagnosis and advise on the best creams. Seek review if the rash changes, becomes painful, itchy, or inflamed, or if there is any doubt about the diagnosis. Persistent or severe cases benefit from dermatology advice.

\textbf{Call 000 (or local emergency services) for:}
\begin{itemize}
    \item Rapidly spreading redness, pain, or swelling, which may indicate infection
    \item Fever with a skin rash
    \item Signs of a severe allergic reaction
    \item Any skin change that is bleeding, growing, or changing rapidly
\end{itemize}

\section*{Diagnosis and Investigations}
\begin{itemize}
    \item \textbf{History and examination:} the characteristic appearance and sites confirm the diagnosis
    \item \textbf{Assessment of associated conditions:} eczema, dry skin, and asthma
    \item \textbf{Review of treatment response:} which products help
    \item \textbf{Skin biopsy:} rarely needed, for uncertain or unusual cases
\end{itemize}

\section*{Management}
\begin{itemize}
    \item \textbf{Regular moisturising:} emollients applied daily, especially after bathing
    \item \textbf{Keratolytic creams:} products with urea, lactic acid, or salicylic acid to soften the plugs
    \item \textbf{Gentle exfoliation:} soft cloths or mild scrubs, without harsh rubbing
    \item \textbf{Avoiding irritants:} harsh soaps, hot water, and tight clothing
    \item \textbf{Bathing practices:} lukewarm showers and patting dry
    \item \textbf{Retinoid creams:} for resistant or severe cases, under medical advice
    \item \textbf{Humidifiers:} in dry environments
    \item \textbf{Realistic expectations:} the condition improves but is usually not eliminated
\end{itemize}

\section*{Complications}
\begin{itemize}
    \item Dry, cracked skin and irritation from aggressive treatment
    \item Redness and inflammation if harsh exfoliation is used
    \item Secondary infection of irritated skin
    \item Cosmetic concerns and effects on confidence
    \item The condition itself is harmless and has no serious complications
\end{itemize}

\section*{Prognosis and Outcomes}
Keratosis pilaris is a chronic but benign condition. Most people see significant improvement with regular moisturising and keratolytic creams, and the bumps flatten and smooth over weeks to months. It often improves with age, and although it can persist, it rarely interferes with life. There is no cure, and treatment is aimed at control, but for most people the condition is easily managed and of cosmetic concern only.

\section*{Prevention and Self-Care}
\begin{itemize}
    \item Moisturise daily, especially after showering
    \item Use lukewarm water and gentle, fragrance-free cleansers
    \item Apply creams containing urea or lactic acid as directed
    \item Exfoliate gently, without scrubbing hard
    \item Pat the skin dry rather than rubbing
    \item Wear loose clothing and avoid friction
    \item Use a humidifier in dry climates
    \item Be patient: improvement takes weeks
    \item See a doctor if the rash changes or causes concern
\end{itemize}

\section*{Sources and Further Reading}
The following reputable sources provide further information for healthcare students and patients seeking reliable, current advice about keratosis pilaris.
\begin{itemize}
    \item Healthdirect Australia. \textit{Keratosis pilaris.} https://www.healthdirect.gov.au/
    \item Australasian College of Dermatologists. \textit{Keratosis pilaris.} https://www.dermcoll.edu.au/
    \item Eczema Association of Australasia. \textit{Dry skin and keratosis pilaris.} https://www.eczema.org.au/
    \item DermNet NZ. \textit{Keratosis pilaris.} https://dermnetnz.org/
    \item NHS. \textit{Keratosis pilaris.} https://www.nhs.uk/conditions/keratosis-pilaris/
    \item Mayo Clinic. \textit{Keratosis pilaris.} https://www.mayoclinic.org/
\end{itemize}
